% Owner: Youssef. Source: docs/related-work.md, which records how well each claim was verified.
% Every \checkcite marks a claim not yet read in the cited source itself; read it and remove the
% marker, or cut the claim, before submission.
\section{Related Work}\label{sec:related}

\paragraph{Pre-execution command guards.}
The closest published work is CARE~\cite{care}, which verifies each shell command before the
host shell runs it. It canonicalises the command with a bounded, deterministic rewrite that
unwraps wrapper shells and normalises quoting, then scores it in four static layers (structural
indicators, a risk class per command head, the sensitivity of the paths it names, and a catalogue
of regular-expression rules weighted by provenance) and consults a language-model judge only for
commands whose score falls in an uncertain band. It reports 85.64\% F1 at a 0.91\%
false-positive rate\checkcite{CARE results table}. Two of its findings shape our design. On one
attack set, the static profile alone reduced realised harm more than the full system with the
judge, apparently because the judge moved dangerous borderline cases to allow%
\checkcite{CARE, RedCode-gen comparison}; our classifier is therefore allowed only to tighten a
verdict. And CARE states that its mediation is stateless, each decision made over a single
command in isolation, naming trajectory-level reasoning as future work. Our design differs from
CARE in three respects: it compiles to effects over a grammar we specify rather than scoring a
normalised string, it resolves paths against session state, and it executes its own canonical
output rather than the agent's command.

The open-source guards share one trait that matters here: none writes its own shell grammar.
CARE parses with bashlex, bashguard~\cite{bashguard} with tree-sitter-bash, and others with
unnamed parsers whose coverage their authors describe only as finite. bashguard separates
evidence from decision, emitting findings with an action type that a policy then turns into a
verdict, and it ratchets: a project configuration may tighten its verdicts but never relax them.
cc-bash-guard~\cite{cc-bash-guard} has the most developed policy language of the group, with
deny, ask and allow buckets, an explicit \emph{abstain} when no rule has an opinion, and a test
block required on every rule that a verifier runs before the policy is accepted. We adopt the
ratchet and the per-rule tests. What none of these tools does is publish the language it
accepts, treat everything outside that language as a refusal, or carry the quoting of each word
through to the decision instead of normalising it away.

\paragraph{Why string filters fail.}
The GuardFall disclosure~\cite{guardfall} attributes the failure of agents' command filters to
inspecting a string before the shell expands it, which is the parser-differential problem in
its plainest form. CmdNeedle~\cite{cmdneedle} shows the same fragility from the other side, using
a language model to propose bypasses and validating each by its side effects in a sandbox, a
method we reuse to build our bypass corpus. \emph{Don't Let AI Agents YOLO Your
Files}~\cite{yolo} names the \emph{shell loophole}: commands run through the shell escape the
policies an agent harness applies to its own file tools, and the paper argues for filtering on
filesystem effects rather than on command strings, which is the premise of our design.
Red-teaming work~\cite{redteam-agentic} argues that enforcement must be non-bypassable by the
agent it constrains, which motivates our rule protecting the gate's own files.

\paragraph{Agent behaviour after a refusal.}
The survey of Dantas et al.~\cite{survey-safe-agents} frames agent safety as specification,
verification and enforcement, and lists the verifier tax among its priority open problems.
\emph{Permission Denied}~\cite{permission-denied} measures post-denial behaviour directly,
evaluating twelve coding agents under nested system-level restrictions and finding that
harnesses which retry aggressively waste effort on denied actions while those that change
strategy keep their success rate\checkcite{beyond the abstract, from secondary coverage}. Its
agents meet raw operating-system errors. A gate that refuses before execution can instead explain
the refusal, and whether the content of that explanation changes what the agent does next appears
to be untested; it is one of our experiments. AgentGuard~\cite{agentguard-trajectories}, which
shares our name by coincidence, learns behavioural rules from failed agent trajectories and
injects them into the agent's context, reporting over-refusal as its main limitation. It
\emph{advises} the agent where we \emph{enforce} at the point of execution, and advice can be
ignored. ClawGuard~\cite{clawguard} addresses indirect prompt injection in tool-using agents at
runtime and is adjacent rather than overlapping.

\paragraph{The shell as a language.}
Our grammar work is not novel in itself, and the shell literature is where the comparison
belongs. Smoosh~\cite{smoosh} gives an executable formal semantics of the POSIX shell, tested
against the standard's conformance suite, and Morbig~\cite{morbig} is a static parser for POSIX
shell written for fidelity to the specification and used to analyse Debian maintainer scripts
at scale. Both are candidate oracles for our differential tests, alongside bash itself. PaSh
\cite{pash} annotates commands with their inputs and outputs in order to parallelise scripts,
which is prior art for our table of command effects. Shill~\cite{shill} lets a script declare the
capabilities it needs, prior art for policies over what a command may touch. And try~\cite{try}
runs a command in an overlay and shows its filesystem changes before committing them, the
dynamic counterpart of predicting effects statically and a natural baseline. What we take from
this literature is the discipline of specifying the language; what we add is applying it at the
boundary between an agent and the shell, where everything outside the specification is a
refusal.
